\section{Taylor coefficients give a stronger all-copy comparison}

\subsection*{Finding M-013: a determinant jet map with surjective differential}
Let \(r\geq2\), put \(R=r-1\), and choose \(\ell\geq0\).
Define
\[
 b=2\ell+2,\qquad h=\lceil\sqrt R\rceil,\qquad
 J_b(r)=2h\binom{R+\ell}{\ell}.
\]
If \(n\geq J_b(r)\), the map from \(n\)-by-\(n\)
linear matrix pencils in \(x_1,z_1,\ldots,z_R\) to the
Taylor coefficients through degree \(b\) of
\(\det A(1,z)\) is dominant. More precisely, its differential
is surjective at an explicitly constructed pencil over \(\mathbb Z\).
Consequently, for every outer degree \(d\) and every
\(\lambda\vdash nd\) with \(\ell(\lambda)\leq r\)
and \(\lambda_2\leq b\),
\begin{equation}\label{eq:M013K}
                         K_n(\lambda)=a_\lambda(d[n]).
\end{equation}
As usual one-row shapes have second row zero. We only discuss
types fitting in the ambient \(n^2\) variables; the displayed
size bound ensures \(r\leq n^2\).
Also \(J_b(r)\geq b\), since
\(\binom{R+\ell}{\ell}\geq\ell+1\) and \(h\geq1\),
so the stated jet order never exceeds the determinant degree.

The construction permits other Taylor coefficients to vary.
It does not represent every polynomial in the truncated space
as an exact padded determinant. Dominance of its detected
coefficients is sufficient to preserve every highest-weight copy.

\paragraph{The acyclic pencil.}
Choose an injection from the \(R\) variable labels into pairs
\((a,c)\in\{1,\ldots,h\}^2\). Denote the corresponding
variable by \(z_{ac}\); unused pairs have label zero.
For each \(a\), take a copy of the monomial tree on
\(\alpha\in\mathbb Z_{\geq0}^R\), \(|\alpha|\leq\ell\),
oriented toward its root. For each \(c\), take a separate
copy oriented away from its root. The parent of a nonzero
multi-index removes its last nonzero coordinate; label that
edge by the corresponding variable, with positive sign.
Connect inward root \(a\) to outward root \(c\) by the
edge labeled \(z_{ac}\) when this is nonzero. No roots
are identified, and no other edges are present.

There are \(J_b(r)\) vertices. The graph is acyclic:
paths run inward, cross at most one bridge, and then run
outward. Let \(B(z)\) be its adjacency matrix. After a
topological ordering it is strictly triangular, hence nilpotent,
and
\[
 A_0(x_1,z)=x_1 I-B(z),\qquad
 \det A_0=x_1^{J_b(r)},\qquad
 A_0(1,z)^{-1}=I+B(z)+B(z)^2+\cdots.
\]
Only finitely many powers occur. Between inward vertex
\((\alpha,a)\) and outward vertex \((\beta,c)\) there
is exactly one path if the bridge is present. Its product is
\[
                         z^\alpha z_{ac}z^\beta.
                                                               \tag{M-013.1}
\]
Thus the corresponding inverse entry is exactly this monomial.
Every monomial of degree between one and \(2\ell+1\)
has this form: choose one of its variables for the bridge
and split its remaining multiset between \(\alpha,\beta\)
of degrees at most \(\ell\). The diagonal inverse entries
also give the constant monomial one. It follows that the
inverse entries span every polynomial of degree at most
\(b-1=2\ell+1\). For \(n>J_b(r)\), adjoin an
\(x_1 I\) block; all these assertions remain true on the
original block after setting \(x_1=1\).

\begin{proof}[Proof of jet dominance]
Let \(\mathcal P\) be the affine space of all linear matrix
pencils of size \(n\) in these \(r\) variables, and let
\(\mathcal J_b\) be the vector space with basis
\(z^\gamma\), \(|\gamma|\leq b\). The polynomial map
\[
 \Phi:\mathcal P\longrightarrow\mathcal J_b,
 \qquad A\longmapsto\det A(1,z)\pmod{(z_1,\ldots,z_R)^{b+1}}
\]
records exactly the desired coefficient jet. A variation of
one matrix entry by a linear form \(H\) has differential
\[
 d(\det)_{A_0}(H E_{ij})
       =H(1,z)\,\det A_0(1,z)\,(A_0(1,z)^{-1})_{ji}
       =H(1,z)(A_0(1,z)^{-1})_{ji}.
\]
This is also the cofactor derivative identity, so no analytic
matrix inverse is required. Taking \(H=x_1\) and using
(M-013.1) yields every monomial through degree \(b-1\).
Taking \(H=z_t\) yields every degree-\(b\) monomial as
well. The differential of \(\Phi\) at \(A_0\) is
therefore surjective onto \(\mathcal J_b\).

Here is an explicit algebraic justification of dominance.
Choose \(\dim\mathcal J_b\) source directions on which
this differential is an isomorphism and restrict \(\Phi\)
to the resulting affine subspace through \(A_0\). Translate
source and target origins and make an invertible linear change
of target coordinates. The restricted polynomial map then
has the form \(F(t)=t+\text{terms of degree at least two}\).
If a nonzero polynomial \(P\) vanished on its image, the
lowest nonzero homogeneous part of \(P(F(t))\) would be
the same as that of \(P(t)\), a contradiction. Thus
pullback on polynomial functions is injective. This proves
that \(\Phi\) is dominant, and completes the first claim.

In fact the chosen differential directions can be individual
matrix-entry coefficient coordinates, each mapping to one distinct
monomial with coefficient one. The corresponding Jacobian minor
is the identity over \(\mathbb Z\). The translated restricted map
already has \(F(t)=t+O(t^2)\) integrally. Its lowest-term argument
therefore proves injectivity of its polynomial pullback after
every commutative base change. This assertion does not require
an algebraic section or a flatness claim; the representation-theoretic
conclusion below continues to be asserted in characteristic zero.
\end{proof}

\begin{proof}[Proof that every ambient copy survives]
Take a highest-weight vector \(w\) of weight \(\lambda\)
in the degree-\(d\) determinant ideal. Nonnegative polynomial
weights show that it depends only on the first \(r\)
coordinates. The \(j=2\) case of F-016 says, more strongly,
that every coefficient generator used by \(w\) has total
degree at most \(\lambda_2\leq b\) outside the first
coordinate. Consequently \(w\) is a polynomial in precisely
the coefficients recorded by \(\mathcal J_b\).

Every determinant of a linear pencil in these variables is
a linear endomorphism substitution into \(\det_n\), and
belongs to \(\Omega_n\). Thus \(w\circ\Phi=0\).
The injectivity of pullback just proved gives \(w=0\).
This applies to every highest-weight vector, not merely to
one chosen copy. The complete ideal multiplicity space at
\(\lambda\) is zero, which is \eqref{eq:M013K}.
Types of zero ambient multiplicity cause no exception.
\end{proof}

\paragraph{Uniform consequence in every power range.}
For \(m\geq2\), every integer \(c\geq6\), and
\(n\geq m^c\), an actual internally padded permanent
multiplicity obstruction must satisfy
\begin{equation}\label{eq:M013uniform}
                              \lambda_2\geq c+2.
\end{equation}
In particular, the BIP range now requires \(\lambda_2\geq27\),
which strengthens the earlier constant twenty-six.

\begin{proof}
If \(q^m_\lambda>0\), then \(r=\ell(\lambda)\leq m^2\).
One-row shapes have full determinant multiplicity, so assume
\(r\geq2\). For every \(\ell\geq3\),
\[
 \frac{2\binom{m^2+\ell-1}{\ell}}{m^{2\ell}}
 \leq\frac{2\binom{\ell+3}{3}}{4^\ell}
 \leq\frac58<1.                                      \tag{M-013.2}
\]
The first inequality follows from the product formula:
each factor \(1+i/m^2\), \(0\leq i<\ell\), decreases
with \(m\), so its largest allowed value occurs at \(m=2\).
The second inequality is equality at \(\ell=3\); the
ratio on increasing \(\ell\) is
\((\ell+4)/(4(\ell+1))<1\).

If \(c=2\ell+1\) is odd, then \(\ell\geq3\),
\(R=r-1\leq m^2-1\), and \(h=\lceil\sqrt R\rceil\leq m\).
For the even jet order \(b=c+1=2\ell+2\),
\[
 J_b(r)\leq2m\binom{m^2+\ell-1}{\ell}
                      \leq\tfrac58m^{2\ell+1}<n.
\]
Thus every type with \(\lambda_2\leq c+1\) has
\(K_n=a_\lambda\) by the jet theorem.

If \(c=2\ell\) is even, use the exact degree-split
construction O-012.a at base degree \(c+1=2\ell+1\).
Its size, for \(r\) variables, is
\[
 2\binom{r-1+\ell}{\ell}-1
       \leq\tfrac58m^{2\ell}-1<n.
\]
Monotonicity in the base degree and F-014 therefore give
\(K_n=a_\lambda\) whenever \(\lambda_2\leq c+1\)
in this case too. In either parity \(q^m_\lambda\leq a_\lambda\)
excludes the alleged obstruction, proving \eqref{eq:M013uniform}.
\end{proof}

\paragraph{Relation to the staircase.}
Combining this result with F-017, the BIP range requires the
complete staircase
\[
                    \lambda_{2j}\geq28-j
                              \qquad(1\leq j\leq26).
\]
The new endpoint is \(\lambda_2\geq27\); the other
twenty-five inequalities are the mixed graph theorem.
The jet argument has changed the required geometric input:
a dominant map to the coefficients detected by a section is
enough. Exact containment of the entire truncated polynomial
family is a stronger condition and is not claimed here.

\paragraph{Grothendieck dictionary.}
The truncated algebra
\(\mathbb Z[z_1,\ldots,z_R]/(z_1,\ldots,z_R)^{b+1}\)
is the explicit finite infinitesimal neighborhood used to
record the coefficients. The graph, pencil and cofactor maps
are defined over \(\mathbb Z\). Their monomial gradings
are torus comodules as in SGA~3, I, Proposition~4.7.3,
pp.~35--36. The flag-minor support is the Schur/exterior
functor calculation of SGA~6, VI, Theorem~3.3, pp.~392--393.
The target jet module is a finite locally free coefficient
module, and the unit Jacobian minor makes its polynomial coordinate
ring inject into that of the pencil parameter scheme after arbitrary
base change. The consequent characteristic-zero vanishing of the ideal
comodule uses SGA~3, I, Proposition~4.7.2, pp.~34--35.
This is an injection of actual coordinate functions, which
contains more information than equality of virtual classes
in \(K_0\). The coherent coefficient quotients fit the
parameter-space framework of FGA, Bourbaki~221, following
Proposition~3.8, pp.~265--266. No unproved positivity of a
motivic trace or integral irreducible splitting is used.
